\SourceLeaf{427}%
das Element $\delta\varrho$, so erhält man
\[
\delta x=X\delta\varrho,\qquad\delta y=Y\delta\varrho,\qquad\delta z=Z\delta\varrho
\]
und
\[
\delta V=P\delta x+Q\delta y+R\delta z=\mu\delta\varrho.
\]
Man sieht also, wie das Vorzeichen von $\mu$ von der Vorzeichenänderung des Werts von $V$ beim Durchgang von der unteren zur oberen Seite abhängt.

\Needspace{5\baselineskip}
\begin{center}9.\end{center}
\noindent\hspace*{1.2em}Man schneide die krumme Fläche mit einer Ebene in dem Punkt, auf den sich unsere Bezeichnungen beziehen. So entsteht eine ebene Kurve, von der $ds\mathord{\hbox{\begin{tikzpicture}[x=.1pt,y=-.1pt,baseline=-7pt]
\path[use as bounding box] (0,0) rectangle (17,101);
\path[fill=black] (14,69) -- (12,69) -- (12,70) -- (9,72) -- (9,74) -- (3,80) -- (2,82) -- (2,85) -- (0,87) -- (0,88) -- (2,90) -- (2,96) -- (3,97) -- (3,99) -- (5,100) -- (11,100) -- (14,97) -- (14,88) -- (13,87) -- (13,84) -- (14,83) -- (14,78) -- (16,76) -- (16,75) -- (14,73) -- cycle (6,3) -- (5,4) -- (5,12) -- (6,14) -- (9,14) -- (14,9) -- (14,6) -- (11,3) -- cycle;
\end{tikzpicture}
}}$ ein Element sei, für das wir die obigen Bezeichnungen beibehalten. Als obere Seite der Ebene betrachten wir diejenige, auf der die Normale zur krummen Fläche liegt. Auf der Ebene errichten wir eine Normale, deren Richtung durch den Punkt $\mathfrak{L}$ der Hilfskugel bezeichnet werde. Beim Fortschreiten entlang der Kurve ändern also $\lambda$ und $L$ ihre Lage, während $\mathfrak{L}$ konstant bleibt und stets $\lambda L$ und $\lambda\mathfrak{L}=90^{0}$ gelten. $\lambda$ beschreibt daher den Großkreis, der $\mathfrak{L}$ als einen seiner Pole hat. Das Element dieses Großkreises ist
\[
=\frac{ds}{r}
\]
wenn $r$ den Krümmungsradius der Kurve bezeichnet. Bezeichnet man wiederum die Richtung eben dieses Elements auf der Hilfskugel mit $\lambda'$, so liegt $\lambda'$ offenbar auf demselben Großkreis und ist sowohl von $\lambda$ als auch von $\mathfrak{L}$ um $90^{0}$ entfernt. Setzt man nun noch
\[
\cos\lambda'(1)=\xi',\qquad\cos\lambda'(2)=\eta',\qquad\cos\lambda'(3)=\zeta',
\]
so gilt
\[
d\xi=\xi'\frac{ds}{r},\qquad d\eta=\eta'\frac{ds}{r},\qquad d\zeta=\zeta'\frac{ds}{r}
\]
da $\xi,\eta,\zeta$ tatsächlich die Koordinaten des Punktes $\lambda$ relativ zum Kugelmittelpunkt sind.
